Com crescimento exponencial de $2\%$ ao ano a partir de 1960,
\[
P(t)=3(1.02)^{t-1960}\ \text{bilhoes}.
\]
Em 2060,
\[
P(2060)=3(1.02)^{100}\approx 21.87\ \text{bilhoes}.
\]

Na escala pedida, $10$ bilhoes correspondem a $1.50$ polegadas. Logo, a altura da ordenada em 2060 seria
\[
\frac{21.87}{10}\times 1.50\approx 3.28\ \text{pol}.
\]

Alguns pontos intermediarios para o tracado sao:
\[
P(1980)\approx 4.46,
\quad
P(2000)\approx 6.62,
\quad
P(2020)\approx 9.83,
\quad
P(2040)\approx 14.60,
\quad
P(2060)\approx 21.87.
\]

O grafico cresce de modo suave, mas o valor de quase $22$ bilhoes em 2060 ja parece agressivo quando comparado com a experiencia historica real. Portanto, como projecao exponencial pura, a curva e matematicamente simples e facil de tracar, mas nao parece especialmente "razoavel" como previsao demografica de longo prazo.